% Options for packages loaded elsewhere
\PassOptionsToPackage{unicode}{hyperref}
\PassOptionsToPackage{hyphens}{url}
\PassOptionsToPackage{dvipsnames,svgnames,x11names}{xcolor}
%
\documentclass[
  11pt,
  a4paper,
]{article}
\usepackage{amsmath,amssymb}
\usepackage{iftex}
\ifPDFTeX
  \usepackage[T1]{fontenc}
  \usepackage[utf8]{inputenc}
  \usepackage{textcomp} % provide euro and other symbols
\else % if luatex or xetex
  \usepackage{unicode-math} % this also loads fontspec
  \defaultfontfeatures{Scale=MatchLowercase}
  \defaultfontfeatures[\rmfamily]{Ligatures=TeX,Scale=1}
\fi
\usepackage{lmodern}
\ifPDFTeX\else
  % xetex/luatex font selection
\fi
% Use upquote if available, for straight quotes in verbatim environments
\IfFileExists{upquote.sty}{\usepackage{upquote}}{}
\IfFileExists{microtype.sty}{% use microtype if available
  \usepackage[]{microtype}
  \UseMicrotypeSet[protrusion]{basicmath} % disable protrusion for tt fonts
}{}
\makeatletter
\@ifundefined{KOMAClassName}{% if non-KOMA class
  \IfFileExists{parskip.sty}{%
    \usepackage{parskip}
  }{% else
    \setlength{\parindent}{0pt}
    \setlength{\parskip}{6pt plus 2pt minus 1pt}}
}{% if KOMA class
  \KOMAoptions{parskip=half}}
\makeatother
\usepackage{xcolor}
\usepackage[margin=25mm]{geometry}
\usepackage{longtable,booktabs,array}
\usepackage{calc} % for calculating minipage widths
% Correct order of tables after \paragraph or \subparagraph
\usepackage{etoolbox}
\makeatletter
\patchcmd\longtable{\par}{\if@noskipsec\mbox{}\fi\par}{}{}
\makeatother
% Allow footnotes in longtable head/foot
\IfFileExists{footnotehyper.sty}{\usepackage{footnotehyper}}{\usepackage{footnote}}
\makesavenoteenv{longtable}
\setlength{\emergencystretch}{3em} % prevent overfull lines
\providecommand{\tightlist}{%
  \setlength{\itemsep}{0pt}\setlength{\parskip}{0pt}}
\setcounter{secnumdepth}{-\maxdimen} % remove section numbering
\usepackage{microtype}
\usepackage{xurl}
\usepackage{fancyhdr}
\pagestyle{fancy}\fancyhf{}
\fancyhead[L]{\small Evidence Press — Perron optimisation}
\fancyhead[R]{\small v1.0}
\fancyfoot[C]{\thepage}
\setlength{\headheight}{14pt}
\setlength{\emergencystretch}{2em}
\ifLuaTeX
  \usepackage{selnolig}  % disable illegal ligatures
\fi
\usepackage{bookmark}
\IfFileExists{xurl.sty}{\usepackage{xurl}}{} % add URL line breaks if available
\urlstyle{same}
\hypersetup{
  pdftitle={Response to the review},
  colorlinks=true,
  linkcolor={Maroon},
  filecolor={Maroon},
  citecolor={Blue},
  urlcolor={Blue},
  pdfcreator={LaTeX via pandoc}}

\title{Response to the review}
\usepackage{etoolbox}
\makeatletter
\providecommand{\subtitle}[1]{% add subtitle to \maketitle
  \apptocmd{\@title}{\par {\large #1 \par}}{}{}
}
\makeatother
\subtitle{Exact Perron optimisation with factorised pair interactions}
\author{}
\date{Version 1.0 · 28 September 2026}

\begin{document}
\maketitle

\subsection{Scope of the revision}\label{scope-of-the-revision}

The review of \emph{Constructive resolvent ordering and weighted Perron
minima}, research draft 0.1, recommended major revision with a
favourable mathematical assessment and a provisional 3-star judgement.
The supplied review and checking archive are preserved unchanged. No
external specialist identity is inferred from them.

Version 1.0 is a complete standalone article centred on the
heterogeneous minimum. It repairs the endpoint and domain statements,
gives explicit predecessor comparisons, and presents the stability
coefficient with its actual limitations. It also supplies new structural
and robust-design results. \textbf{Those additions have written proofs
and exact checks, but are not covered by the supplied referee's
favourable proof assessment.} No revised star rating is claimed.

\subsection{Required revisions}\label{required-revisions}

\begin{longtable}[]{@{}
  >{\raggedright\arraybackslash}p{(\columnwidth - 2\tabcolsep) * \real{0.5000}}
  >{\raggedright\arraybackslash}p{(\columnwidth - 2\tabcolsep) * \real{0.5000}}@{}}
\toprule\noalign{}
\begin{minipage}[b]{\linewidth}\raggedright
Review request
\end{minipage} & \begin{minipage}[b]{\linewidth}\raggedright
Action and location
\end{minipage} \\
\midrule\noalign{}
\endhead
\bottomrule\noalign{}
\endlastfoot
4.1: theorem-level novelty comparison & Section 1.1 distinguishes
rank-one Levinger theory, single arc reversals, independent row
permutations, diagonal optimisation, classical tetrads, and the
content-identified predecessor. Full-text limitations are explicit. \\
4.1: stable predecessor identity & The bibliography gives the
predecessor's complete title, version, date and full archive SHA-256.
The exact archive is nested in the retained reviewed input; its content
identity is independently checked. No DOI has been invented. \\
4.2: coefficient/slack diagnostics & Section 6.1 reproduces the
referee's two six-point diagnostics, explicitly noting that the uniform
bound does not beat the trivial 15-edge bound there. The independent
referee program is replayed. \\
4.2: worked recovery example & An entirely rational parameter choice
yields a four-edit return although a one-edit copy exists. All 720
relabellings are checked. The uniform, exact remainder and
instance-dependent coefficient are separated. \\
4.3: endpoint & Section 5.2 restricts spectral transfer to even
\(n\ge4\). It explicitly excludes the \(n=2,t=1\) endpoint, while
continuous resolvent lemmas retain their valid \(n\ge2\) scope. \\
4.3: matrix domain & Section 5.2 says `skew tournament' where the
numerator sign needs parity. It gives \(S=0\) as a counterexample to
extending that step to the full box. \\
4.4: strongest conclusion first & Conventional abstract and introduction
lead with local-minimum classification, global attainment, fixed-data
cospectrality and the scalar minimum. The complete proof is Section
2. \\
\end{longtable}

The proof-critical recovery code already enforced the correct tournament
domain and \(n\ge4\); the review exposed an overbroad manuscript
statement, not an accepted invalid runtime input. Both guards have
explicit negative regression tests in the publication package.

\subsection{Mathematical
strengthening}\label{mathematical-strengthening}

\textbf{Proposition 3.1 and Theorems 3.2--3.3.} Differentiating the
explicit scalar equation gives positive mass/diagonal sensitivities and
strict decrease with the directional parameter. One fixed transitive
orientation is optimal at every parameter instance. This gives minimax
equality for a compact uncertainty family and identifies the exact
corner for rectangular uncertainty. The robust stability test is a
rational polynomial sign. The restriction to a fixed unknown matrix,
rather than arbitrary switching, is explicit. These are consequences of
the common optimiser, not claims of a new abstract minimax theorem.

\textbf{Proposition 3.4.} A multiplicatively accurate factorisation
supplies a \(\kappa^2\) approximation guarantee for any transitive
order. The factor is not claimed sharp.

\textbf{Theorem 4.1.} Factorisation is necessary and sufficient for
transitive-order cospectrality for every diagonal choice. An exact
four-cycle determinant difference exposes the classical tetrad
obstruction. The classical rank-one equations are attributed to their
factor-analysis context. The theorem does not characterise every
possible no-spurious-minimum family.

\textbf{Proposition 4.2.} Arbitrarily small departures from
factorisation can create strict nonglobal local minima. The proof uses
continuity of strict inward derivatives at transitive vertices and an
exact nonzero constant difference between two characteristic
polynomials. At the rational example \(b_{12}=101/100\), \(t=1/2\), all
24 transitive vertices are certified strict local minima and have two
different Perron values. The larger ones are therefore nonglobal. There
is no claim that the smaller class is the global minimum over that
nonfactorised box.

These additions address the review's significance concern by explaining
both why the weighted principle yields a universal design rule and where
its exact landscape fails. They do not establish historical priority or
a four-star assessment.

\subsection{Smaller corrections}\label{smaller-corrections}

The entrywise discrepancy is now explicitly \(E=2k\), not confused with
edge distance. Matrices with zero diagonal are described as having
positive off-diagonal entries. The one-sided second-derivative argument
is expanded. Arithmetic complexity is separated from bit complexity and
from general interval treatment of an input \(t\). Assurance material is
in Appendix A and the evidence ledgers. The core text no longer requires
the reader to reconstruct conversational provenance.

The title, publication date and local editorial release identifier are
fixed in \texttt{RELEASE.json}. The manuscript remains anonymous. Public
DOI/URL, accountable steward and an effective public reuse licence
remain publisher-controlled fields. A proposed licensing choice is
recorded in \texttt{publisher/LICENCE\_HANDOFF.md}; it is not silently
activated. Thus the review's licensing suggestion is handed off
explicitly, rather than represented as author-approved.

\subsection{Source comparison and remaining access
limits}\label{source-comparison-and-remaining-access-limits}

Both supplied Psarrakos--Tsatsomeros papers remain the full-text basis
for their comparison. Relevant statements of Engel--Sergeev's
row-permutation characterisation and Drton--Sturmfels--Sullivant's
tetrad result were inspected in primary preprints during this revision.
The publisher record for the latter confirms the 2007 issue date and
DOI.

Kirkland's 1995 published body and Drury's complete published proof were
not obtained. Their publisher/author records support the scoped
statements used, not a proof-level novelty clearance. The fixed-trace
diagonal-perturbation paper is likewise abstract-level context. Search
queries, consulted sections and failures are recorded in
\texttt{SOURCE\_LEDGER.md}. No claim is justified merely by a missing
search hit.

\subsection{Evidence and acceptance
criteria}\label{evidence-and-acceptance-criteria}

The reviewed eight-job full replay passed with its 40-file manifest. The
new Fraction checker reconstructs four-point determinants by permutation
expansion. The independent symbolic checker reconstructs characteristic
polynomials and uses rational interval arithmetic on Perron adjugates to
certify all 144 inward-coordinate signs in the fragility example. The
robust example, factorisation reconstruction and non-nearest recovery
example have exact receipts.

The final replay checks file integrity before and after execution. It
reruns the reviewed package and the independently supplied referee
program in temporary copies; historical originals remain unchanged.
Timing is excluded only from the referee result comparison, never from a
mathematical check. Every failed, changed or tolerance-tested quantity
is identified in the resulting receipt. Negative controls recompute
disposable manifests before testing mathematical corruption under
ordinary and optimised Python. The bundle is not proof-assistant
verification, external specialist endorsement, a deployed website
release or an official research-assessment decision.

\end{document}
